\begin{abstract}
Las constelaciones de satélites en Redes No Terrestres (NTN) representan un pilar fundamental para la conectividad global 6G, pero plantean desafíos significativos en sostenibilidad orbital y eficiencia energética. Este artículo presenta un marco integral para minimizar la generación de debris orbital mediante protocolos de disposición post-misión optimizados y técnicas de maniobra autónoma de evitación de colisiones, junto con estrategias de gestión de energía para enlaces de larga duración que incorporan beam hopping, apagado selectivo de payloads y optimización orbit-aware. Se evalúa el impacto en megaconstelaciones LEO mediante simulaciones basadas en modelos 3GPP NTN y métricas de ciclo de vida. Los resultados demuestran reducciones de hasta el 40% en consumo energético y mejoras sustanciales en la densidad orbital sostenible, contribuyendo a la mitigación del síndrome de Kessler mientras se mantiene la calidad de servicio. Se discuten implicaciones regulatorias y direcciones futuras para diseños net-zero en NTN.
\end{abstract}